\uuid{Pe4s}
\exo7id{4066}
\titre{exo7 4066}
\auteur{quercia}
\organisation{exo7}
\datecreate{2010-03-11}
\isIndication{false}
\isCorrection{false}
\chapitre{Equation différentielle}
\sousChapitre{Equations différentielles linéaires}
\cpge{mpsi-mp2i.primitives-equations-differentielles.c, pcsi.techniques-calcul-primitives.c, ptsi.techniques-calcul-primitives.c, mpsi-mp2i.applications-lineaires.a, pcsi.applications-lineaires.a, ptsi.applications-lineaires.a, mpsi-mp2i.applications-lineaires.b, pcsi.applications-lineaires.b, ptsi.applications-lineaires.b, bcpst1.equations-differentielles-lineaires.b, bcpst1.applications-lineaires-matrices, mp-mpi.equations-differentielles-lineaires.b, psi.equations-differentielles, pt.equations-differentielles.a, bcpst2.applications-lineaires-matrices.a}
\module{Analyse}
\niveau{L2}
\difficulte{}

\contenu{
\texte{
Soit $(p,q)\in\R^2$ et $(E)$ l'équation~: $y'' + p y' + q  y = 0$.
On note $S$ l'ensemble des solutions de $(E)$ et $D$ l'application de $S$ dans 
$S$ définie par $D(f) = f'$.
}
\begin{enumerate}
    \item \question{$D$ peut elle être une homothétie ?}
    \item \question{Déterminer les valeurs de $p$ et $q$ pour lequelles $D$ n'est pas un 
    isomorphisme de~$S$.}
    \item \question{Vérifiez que $D\circ D + p D + q\mathrm{id}_S=0$ et montrer qu'il existe des 
    nombres complexes $r_1$ et $r_2$ tels que :
    $(D - r_1 \mathrm{id}_S) \circ (D - r_2 \mathrm{id}_S) = 0$.}
    \item \question{Les applications $D - r_1 \mathrm{id}_S$, $D - r_2 \mathrm{id}_S$
    peuvent-elles être inversibles ?}
\end{enumerate}
}
